\appendix
\section{Supplementary Methods}
\label{app:methods}
\emergencystretch=3em

This appendix gives the settings needed to rebuild the system that the main text describes only in outline. The code, checkpoints and result files are in the supplementary release.

\subsection{Perception Agents}
Both senders were trained with the cross-entropy loss and the Adam optimiser at a learning rate of $3\times10^{-4}$, a batch of 256 windows and at most 30 epochs, stopping when the validation loss had not improved for five epochs. The training set was resampled with class weights proportional to inverse class frequency, with the weight of any class capped at 20 times that of the majority class; uncapped weighting (the fusion class is about one hundredth of the normal class) made validation unstable. Validation used three whole training-split recordings (109, 205 and 223), so that no validation beat shared a patient with training. The encoder used downstream was trained once with random seed 0; ten further seeds were trained to measure its variability (Section~\ref{sec:system}).

The RR branch takes five features per beat computed within a recording: the interval before the beat, the interval after it, the mean of the preceding intervals, and the ratios of the first two to that mean. A single linear layer with a rectifier maps them to 16 dimensions. In the convolutional--recurrent sender these are concatenated to the bidirectional LSTM's output at every time step before the 32-unit context LSTM. The second sender is a one-dimensional residual network: a convolution of width 7 and stride 2 with 32 channels, six residual blocks of two width-5 convolutions with batch normalisation (32, 32, 64, 64, 128 and 128 channels, with stride 2 at the third and fifth), global average pooling to 128 dimensions, concatenation with the 16-dimensional RR embedding, and a linear layer with hyperbolic-tangent activation to the 32-dimensional context vector, followed by the five-class linear head.

\subsection{Prompts and Answer Schema}
Every text arm and the adapter arm share one scaffold: a system prompt, a single class-neutral background paragraph that states the escalation logic without naming the current beat class, a one-sentence note about the window, the event data, and the output instructions. No demonstration examples appear in any multi-event prompt. The system prompt reads: ``You are a clinical decision support assistant reviewing a single ECG beat classification. You will be given structured signal features and fixed background context for the detected beat class. Based on this information only, output a structured urgency tier and a brief justification. Do not invent information not present in the input. This is a research prototype; your output is not sanctioned clinical advice.'' The note reads: ``The event data below contains $N$ consecutive beats from one recording. Apply the rule to each beat and report the single most urgent tier among them.'' The output instructions state the rule exactly (\emph{routine} only for a normal beat; \emph{urgent} if either a run of three or more abnormal beats or a ventricular or fusion beat with confidence above 0.85 holds, each sufficient alone; \emph{priority} otherwise) and ask for a JSON object with three fields: \texttt{urgency\_tier}, a one-sentence \texttt{justification}, and \texttt{referenced\_guideline\_fact}.

The text payloads differ as follows. \emph{A-compact} lists one JSON object per event with four fields: \texttt{label}, \texttt{confidence}, \texttt{consecutive\_abnormal\_beats} and \texttt{signal\_quality\_index}. \emph{A-filtered} replaces the list by one object with \texttt{total\_beats}, \texttt{normal\_beats\_not\_listed} and \texttt{non\_normal\_beats}, a list of the same four-field objects each with its position (1 for the first beat), preceded by a sentence explaining that normal beats are given only as a count. \emph{A-full} uses the event record of the single-event study (class, description, confidence, top three classes, four signal features, signal quality and run length). For the adapter arm the event payload is replaced by the $4N$ virtual tokens and the rest of the scaffold is unchanged. The escalation flag itself is never shown to any arm. The reference tier of an event, and of each training target, comes from that flag.

\subsection{Constrained Decoding}
Decoding is greedy with at most 128 new tokens. A logits processor forces the literal parts of the answer object (the key names and quotation marks, using the tokenizer's own tokenisation), restricts the tier to the three valid words and lets the model's scores choose among them, and allows the two free-text fields any token that contains no double quotation mark, backslash, brace or control character. A field closes when the model emits a token that begins the closing literal, or after 40 tokens, and cannot close while empty. After the closing brace only the end-of-sequence token is allowed. The same processor is used for every arm. For the confirmatory test the text arms stop at the tier token, because stopping there changed no tier on the primary test set.

\subsection{Adapter Initialisation, Data and Selection}
The projection weights are drawn from a normal distribution with standard deviation $0.02/\sqrt{35}$ and the bias is zero. The position embeddings are drawn from a normal distribution with standard deviation 0.02. The scale $\gamma$ is stored as a logarithm and initialised to the median norm of the embeddings the language model produces for 4,096 vocabulary tokens drawn at random (fixed seed), so that virtual tokens start at the scale of real ones. Memory is kept near 10~GB by computing logits only at the target positions and using gradient checkpointing in the frozen model (peak 9.5 to 10.3~GB).

Training windows are drawn from the training split, excluding the three validation recordings, stratified by window tier for $N\in\{1,5,10,20,50\}$ with 25 windows per tier and $N$, the class of each window's most urgent beat taken in round-robin order within the cell. Hard-negative windows are routine windows whose lowest beat confidence is below 0.8, at most 15 per $N$, at most three per recording. Validation windows are drawn in the same way from the three held-out recordings (10 per tier and $N$, at most five per recording, plus up to five hard negatives per $N$). Cells that cannot be filled are left short, which gives the 447 training and 170 validation windows reported in Section~\ref{sec:system}. The validation score is balanced accuracy of the greedy answer (96 new tokens), computed every 111 updates, a quarter of an epoch. Early stopping needs three evaluations without improvement and cannot occur before 894 updates (two epochs), and the best checkpoint is kept. The three adapters were trained with random seeds 101, 202 and 303.

\subsection{Calibration of the Text Arms}
For a text arm the tier is the largest of three logits at the tier token. Calibration adds a constant to the routine logit. The constant is chosen on a grid from $-4$ to $8$ in steps of 0.25 as the value that maximises balanced accuracy on the validation windows (ties resolved to the larger value), separately for each text arm, and is then frozen. The selected values were $-3.75$ for compact text and $8.0$ for filtered text. The second is the upper end of the grid; calibration changed the accuracy of filtered text by about 0.01 or less.

\subsection{Recoverability Decoders}
For each field the decoder is trained on samples from the training split and tested on samples from the test split. Beats are drawn per class (normal, supraventricular, ventricular, fusion), 1,500 per class for training and 500 for testing, and placed at slot 0, 9, 19 or 49 of a 50-event window. For the tokens the features are the $4\times2560$ values of the slot, standardised and reduced to 32 principal components; the reference decoders on the 32-dimensional and 35-dimensional inputs use the standardised inputs only. The multilayer perceptron has one hidden layer of 64 units and is trained for at most 2,000 iterations; the linear decoder is a linear classifier or regressor on the same features. The classification fields are balanced accuracy for the beat label, the urgency tier and the urgent flag, and the regression fields are $R^2$ for heart rate and RR interval. Run status uses a separate sample of 4,358 training beats (450 with a run of three or more) and 1,888 test beats (69 with such a run).
